% !TeX program = xelatex

\documentclass[a4paper, 12pt]{article}

\usepackage{amsmath, mathtools}
\usepackage{blindtext}
\usepackage{fontspec}
\usepackage{layout}
\usepackage[
    verbose,
    paper=a4paper,
    % showframe,
    left=1in,
    right=1in,
    top=36pt,
    textheight=698pt,
    headheight=16pt,
    headsep=20pt, %top+headh+heads=72pt
    marginparwidth=0pt,
    marginparsep=0pt,
    includeheadfoot]{geometry}
\usepackage{fancyhdr}
\usepackage{lastpage}
\usepackage{titlesec}
\usepackage{xcolor}

\mathtoolsset{showonlyrefs}

% Set main font
\setmainfont{Times New Roman}

% Header and Footer
\pagestyle{fancyplain}
\fancyhf{}
\pagenumbering{arabic}
\fancyhead[C]{
    \textbf{Momentum MSCA Programme - Call 2}
}
\fancyfoot[R]{\scriptsize
    page \thepage \,out of\, \pageref{LastPage}}
\renewcommand{\headrulewidth}{0pt}
\renewcommand{\footrulewidth}{0pt}

\titleformat{\section}
    {\normalfont\bfseries}{\thesection.}{1em}{}

% Bibliography setup
\usepackage[
    style=ieee,
    backend=biber,
    url=false,
    doi=false,
    maxnames=4]{biblatex}
\addbibresource{refs.bib}
\defbibheading{bibliography}{\section*{References}}

\def\CC{{C\nolinebreak[4]\hspace{-.05em}\raisebox{.2ex}{\scriptsize\bf ++}}}

%%%%%%%%%%%%%%%%%%%%%%%%%%%%%%%%%%%%%%%%%%%
\begin{document}

\begin{center}
\fbox{\parbox{0.98\textwidth}{\centering
\textbf{Research proposal}
}}
\end{center}

\begin{center}
    \renewcommand{\arraystretch}{1.6}
    \begin{tabular}{|p{0.26\textwidth}|p{0.68\textwidth}|}
        \hline
        Name of the Applicant: &
        Felipe Eduardo Ponce Vanegas \\
        \hline
        Application Title: &
        A Bayesian Approach to Nonlinear System Identification \\
        \hline
    \end{tabular}
\end{center}

\section{Excellence}

Reliable, autonomous, and smooth running machines is the most desirable state for any ma\-nu\-facturing plant.
To achieve, or approach, this utopian scenario it is essential to understand
the response of the machines to external forces so that
unstable operating conditions can be avoided or controlled.
To understand the response of a structure to external forces --- an activity known as System Identification ---
analysis of experimental data and modelling are essential components.

The general objective (GO) of this proposal is
\emph{to provide sound methods and algorithms for the analysis of experimental data
of structures with significant nonlinear effects.}
Here, we focus on structures encountered in Advance Manufacturing by metal removal operations,
like manufacturing of aerospace thin-walled components which
must comply with very high standards of quality and safety.

Unlike many successful data-driven models for image or language processing,
the amount of data available for training of models in Advance Manufacturing is much more limited.
To compensate for such scarcity of data (or relevant information),
physically motivated models are still in active use in the field.

Physical models of structures usually take the form
\begin{equation} \label{eq:mechanical_model}
    M\ddot{x} + C\dot{x} + Kx + G(t, x, \dot{x}) = f(t),
\end{equation}
where \(x\) are general coordinates of displacement from equilibrium,
\(f\) is an external force or excitation,
\(M\), \(C\), and \(K\) are matrices representing the mass, damping, and stiffness, and
\(G\) represents additional linear or nonlinear effects.
Linear models have been successfully applied in the modelling of many structures of practical interest, like machine tools.
However, in recent years, nonlinear effects are getting increasing attention
for, among other things, their relevance in lightweight structures \cite{martinelli24}, or
also for describing the dynamics of robots \cite{bjorn25}.

Nonlinear System Identification is an endless source of challenges.
I plan to approach the subject from a Bayesian point of view.
Bayesian statistics offers a rigorous answer to the question:
to what extent do the collected data support a model?
This is the scientific question we want to answer from our experiments, which
is in clear contrast with the popular Frequentist approach,
where a single model is assumed to be true, and
we have to find it by fitting parameters.
Bayesian statistics is particularly important when decisions must be made with little relevant information, and
we want to include modelling uncertainties into the predictions.


Bayes's theorem states that
\begin{equation} \label{eq:bayesT}
    p(M \,|\, D) \propto p(D \,|\, M)\,p(M),
\end{equation}
which means that the posterior, or conditional probability \(p(M \,|\, D)\) of a model \(M\) given the observed data or evidence \(D\), is
proportional to the likelihood, or conditional probability \(p(D \,|\, M)\) of the data given the model, times
the prior, or probability \(p(M)\) of the model.
The posterior \(p(M \,|\, D)\) is the central object to make predictions, and
Bayes's theorem is the tool to compute it.

To make Bayesian inference a viable alternative,
in agreement with the general objective (GO),
my first objective (O.1) is
\emph{to explore and apply methods to speed up the computation of the likelihood.}
This objective is roughly equivalent to efficiently solving \(x\) in \eqref{eq:mechanical_model}.
Two notable barriers that I must overcome are that
(i) the number of variables could be very large, and
(ii) standard solvers may be unacceptably slow.

To address (i),
existing techniques like substructuring, Reduction Order Models, and Nonlinear Normal Modes could be exploited.
On the other front,
standard solvers like Runge--Kutta schemes may be very slow, so
for (ii) other alternatives like Neural Networks offer an interesting alternative.
Solutions to both problems have been proposed, for example, in \cite{conti23}.

Prior and historical knowledge is valuable when there is little data, so
my second objective (O.2) is
\emph{to enable integration of prior knowledge through \(p(M)\).}
The Bayesian theoretical framework elegantly updates information using a previous posterior as a new prior.
Unfortunately,
technical issues make it difficult to realize Bayes's theorem in practice.
Recently, Neural Networks have been proposed as a fast
alternative to compute probability densities \cite{cai24},
so I will explore and further develop these kinds of algorithms.

Incidentally, O.2 naturally leads to researching algorithms for online updating,
which are important for structures whose parameters substantially vary in time, like for robots.
Hence, I set up as third objective (O.3)
\emph{to develop algorithms for online updating.}
The activities to reach O.2 are also included here, however,
I can optionally explore better methods for computing predictions exploiting decentralized and federated learning,
taking into account real computing capabilities in real industrial environments.

The fourth objective (O.4) is
\emph{to validate a methodology with a real experimental system.}

Submitted versions of the articles will be published in the 
Repository of the Academy's Library (REAL), as encouraged by MTA, and also in arXiv.
Depending on the publisher's policies,
the accepted or published versions will be also deposited in REAL;
no funds will be designated for the payment of open access fees.

For the scientific code,
I will give priority to open-source software, and
I will upload the code to public repositories, like GitHub.
Experimental datasets will be uploaded to the repository of the
Hungarian Research Network (HUN-REN).
For more details, see the Data Management Plan.


\section{Impact}

Nonlinear System Identification is a very complex task, and
it may define my whole future academic career.
In the short-term, this project will be an opportunity
to gain more experience and deepen into modern techniques in high demand inside and outside academia,
like: Neural Networks; FEM simulations; data management; and monitoring of industrial processes.
\textcolor{red}{Do you suggest any skills? Social skills? Maybe, I should describe them in more detail.}
\textcolor{red}{Do you think I should take a course on Ansys, or whatever?}

Bayesian statistics is conceptually and computationally more demanding than the standard Frequentist statistics,
so research in this area will also make me stand out among my scientific peers, and
it will give me an advantage in the rigorous analysis of experimental data with robust uncertainty quantification.

The European Commission views Advance Manufacturing as a strategic sector
for maintaining the region's position as a competitive global player \cite{industrialforum}.
In the same way as Modal Analysis is an essential tool today in the manufacturing industry,
the new theoretical frameworks and methodologies developed during this project will
contribute to the creation of new tools for monitoring, control, and diagnosis of machining processes where
nonlinear effects and varying parameters play an important role.

Research outputs will be disseminated through usual scientific channels:
articles in re\-pu\-table peer-reviewed journals, and
posters and talks in scientific events. In particular,
I plan to attend
\textcolor{red}{any well-known conference I should mention as a target event?}

Additionally, the research outputs will be promoted through the webpages of the Department of Applied Mechanics and the Machine Tool Vibration Research Group.
I will also post blog entries to my webpage featuring my research or interesting facts in a more engaging and less formal language for a broader scientific audience.

To grow and strengthen my collaborative network,
I will organize at least one workshop on Nonlinear System Identification and Bayesian statistics to bring together researchers
from both fields.

One product will be an open-source Python package ---
inspired by the principles of NumPy and similar scientific packages ---
where I will collect, organize, and maintain algorithms so that
other researchers can exploit them under a license like MIT or CC-BY.
The algorithms will initially be inspired by my research, but
not constraint to it.
The code will be available in a Git repository, and
the package will be listed in the Python Package Index (PyPI).
A webpage will be dedicated to tutorials and documentation.
For all this, I will take a course on good practices for code development and deployment.

The Python package will not only be intended for the benefit of other researchers, but
also to make new algorithms and paradigms accessible to the industry,
not only companies with big budgets, but also small ones and start-ups.
A successful package would also secure me a consultancy role for companies with relevant use-cases,
which would in turn enrich the content of the package.

\textcolor{red}{Any ideas for ``communication activities for engaging with the broader public'':
i) Target audiences ii) Key messages iii) Tools and channels to be used.
I do not want to commit myself to participate in events for general public (also, I do not speak local people's language).
I do not have Twitter, Facebook, or similar poisons, I avoid posting anything in LinkedIn.
By ``broader public'' I understand people in the industry; any event to sell algorithms?}


\section{Implementation}

As a MSCA fellow,
I will be part of the MTA-BME Lendület Machine Tool Vibration Research Group, led by Zoltán Dombóvári, 
which functions within the Department of Applied Mecha\-nics.
Zoltán Dombóvári is well-known for his contributions to ... \textcolor{red}{you can shine here!}

\begin{center}
    \includegraphics[width=0.9\textwidth]{images/gantt.eps}
\end{center}

Explain risks ...

WP: Dissemination and Transfer Knowledge: Package website. Conferences. Organize workshop.

WP: Management: Data Management Plan. Progress meetings. Secondment.

The main objective of the Machine Tool Vibration Group is to
develop a framework for effective intervention of machining processes based on sensory data.
Among the strategic objectives,
robot dynamics and machining of slender workpieces feature prominently due to
the emergence of nonlinear behavior and time varying parameters,
which perfectly aligns with the objectives of this proposal,
although the approaches are different.

I will bring to the Group a Bayesian perspective which
they have not explored yet.
This new perspective will extend the Group's research portfolio
by endowing them with new tools for the rigorous statistical analysis of experimental data, and
additional theoretical knowledge to include uncertainty in their models.

The Group excels at the construction, simulation, and experimental validation of physical models of machining processes.
However, at times, their work depends on the application of heuristic algorithms or black-boxes
whose inner workings do not lie in their domain of expertise.
On the other hand, as a mathematician,
I rely on their expert knowledge and experience on machining processes.
This creates a perfect synergy because
I can complement them by opening and modifying black-boxes (like Neural Network architectures) tailored
for specific use-cases, while
steering my work toward physically relevant mathematical concepts.

In fact, I have already collaborated with members of the Group in the study of processes with varying parameters,
where they identified an interesting phenomenon in tool dynamics, and
I assisted them by providing rigorous mathematical foundations to their intuitions.

\textcolor{red}{Does BME-MM have access to any computing cluster?}

\textcolor{red}{Does BME-MM have licenses for relevant programs, like: Ansys, Siemens Testlab, Matlab, Abaqus, ...}

\textcolor{red}{About equipment and facilities, what do you suggest I should remark as strengths of MM-BME.}

A secondment is planned at the Royal Institute of Technology (KTH) in Stockholm, Sweden,
under the supervision of Andreas Archenti, who
is an exceptional researcher in precision engineering for
sustainable manufacturing of high-quality products,
which include aerospace components.

During my years in Spain,
I have collaborated with BCAM researchers like Michael Barton in path planning and modelling of metal cutting processes, and
I have participated in projects with Aritz Pérez on intelligent manufacturing.

Also, I was in constant contact with engineers and researchers from the Aeronautics Advanced Manufacturing Center (CFAA), and IDEKO,
both located in the Basque Country, Spain.
These centers have highly qualified researcher on Advance Manufacturing,
who can contribute their knowledge and experimental data.
In particular, IDEKO has a long experience with real industrial problems,
which can inspire new developments in my research.



\printbibliography

\end{document}